% ═══════════════════════════════════════════════════════════════════════════
% §6  WHAT MODULATES ALIGNMENT
% ═══════════════════════════════════════════════════════════════════════════
\section{What Modulates Alignment}
\label{sec:ablation}

\begin{table}[t]
\centering
\small
\setlength{\tabcolsep}{3.5pt}
\renewcommand{\arraystretch}{0.96}
\caption{Per-model alignment ranking. $\checkmark$ = within human LOO 95\,\% CI; $(\checkmark)$ = concentrates on consensus answers beyond human variability. Abst(\%) = abstention rates with (w/) and without (w/o) instruction (\textbf{bold} = best; \underline{underline} = second-best).}
\label{tab:alignment_ranking_7b_ticks}
\begin{tabular}{>{\raggedright\arraybackslash}p{2.6cm}cccccc}
\toprule
\multirow{2}{*}{\textbf{Model}} & \multicolumn{2}{c}{$\djs$} & \multicolumn{2}{c}{$\sHM$} & \multicolumn{2}{c}{Abst(\%)} \\
\cmidrule(lr){2-3}\cmidrule(lr){4-5}\cmidrule(lr){6-7}
 & Yes/No & Count & Other & $\rho$ & w/ $\downarrow$ & w/o $\uparrow$ \\
\midrule
\multicolumn{7}{c}{\textit{VLMs}}\\
\midrule
InternVL3.5 & \textcolor[HTML]{2E7D32}{($\checkmark$)} & - & \textcolor[HTML]{2E7D32}{$\checkmark$} & \textcolor[HTML]{2E7D32}{($\checkmark$)} & 2.1 & 0.0 \\
LLaVA-1.5 & - & - & - & \textcolor[HTML]{2E7D32}{$\checkmark$} & 3.5 & 3.0 \\
LLaVA-Mistral & - & - & \textcolor[HTML]{2E7D32}{$\checkmark$} & \textcolor[HTML]{2E7D32}{$\checkmark$} & 1.8 & 2.6 \\
LLaVA-Vicuna & - & - & - & \textcolor[HTML]{2E7D32}{$\checkmark$} & \underline{0.9} & 2.7 \\
Qwen3-VL & \textcolor[HTML]{2E7D32}{$\checkmark$} & - & \textcolor[HTML]{2E7D32}{$\checkmark$} & \textcolor[HTML]{2E7D32}{$\checkmark$} & 3.9 & 5.3 \\
\midrule
\multicolumn{7}{c}{\textit{Backbone Decoders}}\\
\midrule
InternVL3.5 & \textcolor[HTML]{2E7D32}{($\checkmark$)} & - & \textcolor[HTML]{2E7D32}{$\checkmark$} & \textcolor[HTML]{2E7D32}{$\checkmark$} & 9.5 & \textbf{17.1} \\
LLaVA-1.5 & \textcolor[HTML]{2E7D32}{($\checkmark$)} & - & \textcolor[HTML]{2E7D32}{$\checkmark$} & \textcolor[HTML]{2E7D32}{$\checkmark$} & \textbf{0.6} & 0.0 \\
LLaVA-Mistral & \textcolor[HTML]{2E7D32}{$\checkmark$} & - & \textcolor[HTML]{2E7D32}{$\checkmark$} & \textcolor[HTML]{2E7D32}{$\checkmark$} & 2.1 & 0.3 \\
LLaVA-Vicuna & \textcolor[HTML]{2E7D32}{($\checkmark$)} & - & \textcolor[HTML]{2E7D32}{$\checkmark$} & \textcolor[HTML]{2E7D32}{$\checkmark$} & \textbf{0.6} & 0.3 \\
Qwen3-VL & \textcolor[HTML]{2E7D32}{($\checkmark$)} & - & \textcolor[HTML]{2E7D32}{$\checkmark$} & \textcolor[HTML]{2E7D32}{$\checkmark$} & 3.2 & 5.3 \\
\midrule
\multicolumn{7}{c}{\textit{SA-LLMs}}\\
\midrule
Qwen2.5-Instruct & - & - & \textcolor[HTML]{2E7D32}{$\checkmark$} & - & 6.5 & 9.1 \\
Qwen3 & - & - & \textcolor[HTML]{2E7D32}{$\checkmark$} & - & 4.7 & \underline{12.7} \\
Qwen3 (think) & - & - & \textcolor[HTML]{2E7D32}{$\checkmark$} & - & 4.7 & \underline{12.7} \\
Vicuna & \textcolor[HTML]{2E7D32}{($\checkmark$)} & - & - & - & 1.5 & 1.2 \\
\bottomrule
\end{tabular}
\end{table}

\paragraph{Visual encoding.}
Let $P(a\mid q)$ denote the model's answer distribution over answers $a$ conditioned on question $q$.
The LM computes $P(a\mid q)$ from the question alone---shaped by VQA co-training to closely mirror human linguistic priors. The VLM computes $P(a\mid q,v_0)$, where $v_0$ is the vision encoder's output for a blank image; since $v_0\!\neq\!\mathbf{0}$, the encoder acts as a nuisance that shifts the distribution away from $P(a\mid q)$ without adding information.
This is borne out empirically: in four of five matched families except for InternVL, removing the vision encoder \emph{improves} $\sHM$ (Table~\ref{tab:matched_family_alignment_tests}).
Difficulty concordance (Spearman $\pr$) shows a consistent pattern in Table~\ref{tab:alignment_ranking_7b_ticks}: every backbone decoder reaches the human LOO CI ($\checkmark$), and so do the three LLaVA VLMs---though on free-text questions specifically the gap widens substantially ($\pr{=}0.450$--$0.628$ vs.\ $0.820$--$0.850$), where difficulty tracking is most meaningful and the encoder's disruption most visible.
InternVL reverses this: its VLM exceeds the CI upper bound (highest $\pr$ across all models, no $\checkmark$ because it is \emph{above} the human range), while its LM falls comfortably within it---the encoder in this case sharpens rather than blurs the model's sensitivity to question difficulty. % by $+0.017$--$+0.069$ (LLaVA families: $+0.041$--$+0.069$; ---showing the encoder distorts both \emph{how well} a model answers and \emph{which} questions it finds easier
Processing a blank image introduces an absence signal that flips LLaVA VLMs from their backbone's yes-bias (42--81\%) to a no-default (15--39\%); Qwen3-VL shows an even stronger reversal (LM: 71--74\%; VLM: 11--41\%), with the 2B variant collapsing to just 11.5\% yes. %---a shift absent in SA-LLMs without an encoder: Qwen3-8B sits at 38.5\% and Vicuna at 57.7\%, reflecting their text-prior baselines rather than any encoder-driven push
InternVL's encoder preserves rather than disrupts this, maintaining the backbone's yes-bias (78.2\% vs.\ 80.8\% for its LM).

To test how much VLMs are affected by uninformative inputs, we probe four blank-image variants (black, gray, noise, white; Appendix~\ref{app:ablation}, Table~\ref{tab:image_ablation_main}): Qwen3-VL shows the widest semantic alignment variation ($\sHM{=}0.490$--$0.628$, span $0.138$), LLaVA-Mistral a narrower range ($0.590$--$0.626$, span $0.036$), while InternVL's semantic alignment is nearly invariant to image type ($0.639$--$0.660$, span $0.021$) and the highest across all conditions---suggesting its language prior dominates regardless of which uninformative image is presented.
Since InternVL and Qwen3-VL share the same Qwen3-8B backbone yet behave oppositely, the prior-disrupting vs.\ prior-preserving split traces to encoder architecture and training recipe \citep{tong2024eyes,mckinzie2024mm1}.
For example, \textit{``What condiments are on it?''} \emph{(Identity; Food; Pron)} elicits \textit{ketchup and mustard} from both InternVL's VLM and LM---matching the human plurality (ketchup: 9 of 40)---while all four other VLMs default to \textit{0} (the encoder triggers count-zero confusion), and LMs recover to \textit{ketchup and mustard} once the encoder is removed. % N ($\sHM{=}0.621$ vs.\ $0.261$ for VLMs).
% Attribute: ``What is the color of the dog's collar?'' (Attribute; Animal) → red from InternVL VLM+BD (human plurality 38%); other VLMs/BDs → black/white; SA-LLMs → black
% \textit{Identity} questions reveal a complementary noun-anchoring effect: \textit{``What type of bread does this appear to be?''} \emph{(Identity; Food; Variant~C)} yields diverse but plausible answers (\textit{wheat}, \textit{white bread}, \textit{sourdough}), yet abstracting one level to \textit{``What type of food does this appear to be?''} \emph{(Variant~B)} causes every model---VLMs, backbone decoders, and SA-LLMs alike---to snap to \textit{pizza}, the dominant food category in VQA corpora, regardless of architecture.
% Human (C): croissant, sourdough, hamburger bun, scone; Variant B pizza anchor: all models converge

\paragraph{Image pretraining.}
% Comparing matched backbone-decoder--SA-LLM pairs isolates the effect of image co-training in the language component itself, independent of the encoder.
We can see the effects of image co-training by comparing LM to SA-LLM. % , with encoder contribution held fixed.
Two VLMs in our set---Qwen3-VL and InternVL---share the Qwen3-8B language backbone.
Both their LMS outperform the Qwen3-8B by $+0.087$ and $+0.103$ in $\sHM$ (Appendix Table~\ref{tab:matched_family_alignment_tests}).
\textit{Attribute} questions---which ask about perceptual properties such as color, texture, or material---show the largest co-training benefit: LMs consistently outperform SA-LLMs by $+0.09$--$+0.13$ in $\sHM$ across all variants, as SA-LLMs default to corpus-frequent responses (\textit{e.g.}~\textit{black}) where LMs retain visually-grounded distributions.
\textit{Identity} questions, by contrast, show near-zero co-training benefit ($\sHM{\approx}0.55$ for LMs and SA-LLMs alike): text-only pretraining already captures comparable distributional knowledge for type- and position-naming. \textit{``What is this player's position called?''} \emph{(Identity; Sports)} draws diverse sports positions from both LMs (\textit{goalkeeper}, \textit{goalie}) and SA-LLMs (\textit{striker}, \textit{forward})---both mirroring the dispersed human distribution (striker=9, goalkeeper=6 of 40).
% Backbone decoders: goalkeeper/goalie across all families; SA-LLMs: Qwen3-8B=striker, Qwen3-think=striker, Qwen2.5=forward

\paragraph{Extended reasoning.} % Qwen3 vs. Qwen3 (Think) with CoT
With chain-of-thought (CoT), Qwen3 gains small but consistent gains: $\sHM$ improves by $+0.02$--$+0.06$ and accuracy by $+0.01$--$+0.04$ across variants, with both entering the human LOO band; answer confidence also rises with longer chains (Figure~\ref{fig:think_abs_sbert_accuracy_vert}).
However, the benefit is size-dependent and non-monotonic. 
At 1.7B and 32B over 60\% of responses contain empty reasoning traces, while 4B generates longer chains but still fails to improve $\sHM$.
32B shows a marginal $\sHM$ gain on the answering subset ($+0.030$) but $\rho$ decreases ($0.793{\to}0.734$) and accuracy collapses due to high abstention ($46.7\%$ with instruction, $8.3\%$ without).
% At 4B, however, think \emph{hurts} $\sHM$ ($\Delta{=}{-}0.06$) with no accuracy gain, suggesting CoT is counterproductive at that scale.
% At 32B, $\sHM$ improves marginally but accuracy collapses due to abstention; 1.7B and 32B also show inverted instruction compliance (abstention rises \emph{with} instruction: $15.7\%{\to}46.2\%$ and $8.3\%{\to}46.7\%$).
% Difficulty tracking ($\rho$) regresses at 8B ($0.807{\to}0.750$) regardless.
More details are shown in Appendix~\ref{app:thinking}.

\paragraph{Scale.}
% We test multiple sizes for three families: InternVL (1B, 2B, 8B), Qwen3-VL (2B, 4B, 8B, 32B), and Qwen3 SA-LLM (0.6B--32B), with Vicuna and LLaVA-Vicuna compared at 7B vs.\ 13B (Appendix~\ref{app:scale_variant}; Table~\ref{tab:scale_alignment_full}).
Across VLM and LM families, alignment peaks at 7--8B with no consistent improvement beyond that scale (Table~\ref{tab:alignment_ranking_full_scale_ticks}): InternVL peaks at 8B ($\sHM{=}0.620$) but dips at 2B ($0.553$), while Qwen3-VL LM is nearly flat across 2B--32B ($0.593$--$0.607$)---scale enriches neither family's scene-plausibility prior.
For SA-LLMs, scale adds deliberative capacity but cannot substitute for that prior: Qwen3 rises from $\sHM{=}0.315$ at 0.6B to $0.554$ at 32B yet remains below InternVL LM 8B ($0.609$), and the trajectory is non-monotonic.
An indirect scale effect is notable: InternVL~(LM) $\sHM$ on pronominalized questions \emph{decreases} from 1B to 8B ($0.555{\to}0.526$), driven by growing count abstention at 8B---which the vision encoder suppresses, partially explaining the reversed VLM--backbone ordering at that size.
The most salient scale effect is abstention: Qwen3-VL 32B abstains spontaneously on 89\% of responses (vs.\ 3--5\% at 2--8B), with instruction reducing this only partially ($\Delta{=}{-}25.0\%$); full abstention rates across all sizes are reported in Table~\ref{tab:alignment_ranking_full_scale_ticks} and Appendix~\ref{app:abstention}.

\paragraph{Instruction compliance.}
When no instruction is given, humans would abstain and double check on the task. 
With instruction to imagine the scene, human abstention to the questions is $<\!0.4\%$ (49 of 13,560 responses) with our classifier. We treat this unique behavior of humans and treat this extreme switch as an alignment dimension. 
As seen in Table~\ref{tab:alignment_ranking_7b_ticks}, many of our models shows minimal abstention even without instruction -- some even 0\%, and those ones sometimes gain in abstention with instruction, which is a clear reversal of human behavior. 
The most aligned model in this case is InternVL~(LM), showing the highest spontaneous abstention ($17.1\%$ without instruction) and reduction in abstention with instruction ($9.47\%$). Qwen3-VL VLM and LM also show this switch with small abstention reduction under instruction (VLM: $5.3\%{\to}3.9\%$; LM: $5.3\%{\to}3.3\%$), though at much lower absolute levels.
In sum, while humans abstain freely without instruction and comply cleanly with it, no model achieves both: models that abstain do so inconsistently, and models that comply do so regardless of context.